\section{Entrenamiento de redes neuronales}

\subsection{El objetivo de la optimización de la pérdida}

El objetivo de entrenar una red neuronal es encontrar los pesos $\mathbf{W}^*$ que minimizan la pérdida empírica $J(\mathbf{W})$:

$$\mathbf{W}^* = \arg\min_{\mathbf{W}} J(\mathbf{W}) = \arg\min_{\mathbf{W}} \frac{1}{n}\sum_{i=1}^{n}\mathcal{L}\big(f(x^{(i)}; \mathbf{W}), y^{(i)}\big)$$

Cabe señalar que $\mathbf{W}$ incluye todos los pesos y sesgos de \textbf{todas las capas} de la red: $\mathbf{W} = \{\mathbf{W}^{(0)}, \mathbf{W}^{(1)}, \ldots\}$.

\begin{figure}[H]
\centering
\includegraphics[width=0.85\textwidth]{slides-images/slide-057.jpg}
\caption{El landscape de la función de pérdida: el objetivo es encontrar el punto más bajo\protect\footnotemark}
\end{figure}
\footnotetext{Diapositivas, página 57.}

\subsection{Descenso de gradiente}

El \textbf{descenso de gradiente} (Gradient Descent) es el algoritmo central para optimizar la función de pérdida. Su idea básica es sencilla y elegante:

\begin{importantbox}{El algoritmo de descenso de gradiente}
\begin{enumerate}
\item \textbf{Inicializar aleatoriamente} los pesos $\mathbf{W}$
\item \textbf{Calcular el gradiente}: $\frac{\partial J(\mathbf{W})}{\partial \mathbf{W}}$, el cual apunta en la dirección en la que la pérdida crece más rápido
\item \textbf{Actualizar los pesos}: dar un paso pequeño en la dirección \textbf{opuesta} al gradiente
$$\mathbf{W} \leftarrow \mathbf{W} - \eta \frac{\partial J(\mathbf{W})}{\partial \mathbf{W}}$$
\item Repetir los pasos 2--3 hasta la convergencia
\end{enumerate}
Aquí $\eta$ es la \textbf{tasa de aprendizaje} (learning rate), que controla la magnitud de cada actualización.
\end{importantbox}

\begin{figure}[H]
\centering
\includegraphics[width=0.85\textwidth]{slides-images/slide-060.jpg}
\caption{Visualización del descenso de gradiente: un descenso iterativo en la dirección opuesta al gradiente\protect\footnotemark}
\end{figure}
\footnotetext{Diapositivas, página 60.}

Una intuición del gradiente: imagina que estás de pie en una zona de colinas, con los ojos vendados, y quieres llegar al punto más bajo. Lo único que puedes hacer es sentir la pendiente del terreno bajo tus pies (el gradiente) y dar un paso en la dirección de mayor descenso. Al repetir este proceso, terminarás por llegar a alguna depresión del terreno.

\subsection{Retropropagación}

Entonces, ¿cómo calcular de manera eficiente el gradiente $\frac{\partial J}{\partial \mathbf{W}}$? La respuesta es la \textbf{retropropagación} (Backpropagation).

La retropropagación utiliza la \textbf{regla de la cadena} (Chain Rule) para propagar el gradiente hacia atrás, capa por capa, comenzando desde la capa de salida. Para cada peso $w$ de la red, la regla de la cadena permite descomponer el gradiente de la pérdida final respecto a ese peso en el producto de varios gradientes intermedios.

Tomando como ejemplo una red simple de dos capas, el gradiente de la pérdida $J$ respecto al peso $w_1$ de la primera capa puede expresarse como:

$$\frac{\partial J}{\partial w_1} = \frac{\partial J}{\partial \hat{y}} \cdot \frac{\partial \hat{y}}{\partial z_2} \cdot \frac{\partial z_2}{\partial z_1} \cdot \frac{\partial z_1}{\partial w_1}$$

\begin{importantbox}{El núcleo de la retropropagación}
La esencia de la retropropagación es la \textbf{aplicación recursiva de la regla de la cadena}. Se calculan gradientes locales empezando por el extremo de salida y se propagan hacia atrás, capa por capa, multiplicando todos los gradientes locales para obtener el gradiente final. Esto hace posible calcular gradientes de manera eficiente en redes profundas con millones de parámetros.
\end{importantbox}

Los marcos de aprendizaje profundo modernos (TensorFlow, PyTorch) realizan automáticamente el cálculo del gradiente de la retropropagación mediante la \textbf{diferenciación automática} (Automatic Differentiation), sin que el desarrollador tenga que derivarlo a mano.

\begin{lstlisting}[caption={Diferenciación automática en TensorFlow y PyTorch}]
# TensorFlow: GradientTape
import tensorflow as tf
with tf.GradientTape() as tape:
    loss = compute_loss(model, x, y)
grads = tape.gradient(loss, model.trainable_variables)

# PyTorch: autograd
import torch
loss = compute_loss(model, x, y)
loss.backward()
# Gradients stored in param.grad
\end{lstlisting}

\subsection{La importancia de la tasa de aprendizaje}

La tasa de aprendizaje $\eta$ es uno de los hiperparámetros más determinantes al entrenar una red neuronal. Su valor incide directamente en el éxito o el fracaso del entrenamiento:

\begin{table}[H]
\centering
\begin{tabular}{lp{9cm}}
\toprule
\textbf{Tasa de aprendizaje} & \textbf{Efecto} \\
\midrule
Muy baja & La convergencia es extremadamente lenta y puede quedar atrapada en un mínimo local \\
Muy alta & El entrenamiento es inestable, la pérdida oscila repetidamente cerca del mínimo o incluso diverge \\
Adecuada & Converge de manera estable a un buen mínimo \\
\bottomrule
\end{tabular}
\caption{Efecto del valor de la tasa de aprendizaje sobre el entrenamiento}
\end{table}

\begin{warningbox}{Trampas comunes al ajustar la tasa de aprendizaje}
\begin{itemize}
\item Una tasa de aprendizaje demasiado alta no solo provoca oscilaciones, sino que puede hacer que la pérdida crezca cada vez más (diverja), porque cada actualización rebasa el mínimo
\item Una tasa de aprendizaje demasiado baja es segura, pero puede requerir un tiempo de entrenamiento sumamente largo y con facilidad queda atrapada en mínimos locales poco profundos
\item En la práctica se usan con frecuencia métodos de \textbf{tasa de aprendizaje adaptativa} (como Adam, RMSProp), que ajustan automáticamente la tasa de aprendizaje de cada parámetro a partir del historial de sus gradientes
\end{itemize}
\end{warningbox}

\begin{importantbox}{Algoritmos de tasa de aprendizaje adaptativa}
En el entrenamiento real rara vez se usa una tasa de aprendizaje fija. Entre los optimizadores adaptativos más comunes están:
\begin{itemize}
\item \textbf{SGD with Momentum}: usa un promedio móvil exponencial del gradiente para acelerar la convergencia
\item \textbf{Adam} (Adaptive Moment Estimation): combina estimaciones del primer y del segundo momento para ajustar la tasa de aprendizaje de cada parámetro por separado
\item \textbf{RMSProp}: normaliza la tasa de aprendizaje mediante un promedio móvil del cuadrado del gradiente
\end{itemize}
La idea central de estos algoritmos es que, en lugar de usar una misma tasa de aprendizaje fija para todos los parámetros, esta se ajusta de manera dinámica según el historial del gradiente de cada parámetro.
\end{importantbox}

\subsection{Descenso de gradiente por minilotes}

El descenso de gradiente estándar necesita calcular la pérdida y el gradiente sobre \textbf{todo el conjunto de datos} antes de poder realizar una sola actualización de los pesos, lo cual implica un costo de cómputo enorme en conjuntos de datos a gran escala.

El \textbf{descenso de gradiente estocástico} (Stochastic Gradient Descent, SGD) representa el otro extremo: en cada paso se usa \textbf{una sola muestra} para calcular el gradiente y actualizar los pesos; es rápido, pero muy ruidoso.

La solución intermedia que se usa en la práctica es el \textbf{descenso de gradiente por minilotes}: en cada paso se toma al azar un lote pequeño (minilote) de muestras del conjunto de datos para calcular el gradiente:

$$\frac{\partial J(\mathbf{W})}{\partial \mathbf{W}} \approx \frac{1}{B}\sum_{k=1}^{B}\frac{\partial \mathcal{L}(f(x^{(k)}; \mathbf{W}), y^{(k)})}{\partial \mathbf{W}}$$

Aquí $B$ es el tamaño de lote.

\begin{table}[H]
\centering
\begin{tabular}{lp{4.5cm}p{4.5cm}}
\toprule
\textbf{Método} & \textbf{Ventajas} & \textbf{Desventajas} \\
\midrule
GD por lote completo & Gradiente exacto, convergencia estable & Alto costo de cómputo y gran demanda de memoria \\
SGD ($B=1$) & Frecuencia de actualización muy alta & Gradiente muy ruidoso, convergencia inestable \\
SGD por minilotes & Equilibra velocidad y estabilidad, se puede paralelizar & Requiere ajustar el tamaño de lote \\
\bottomrule
\end{tabular}
\caption{Comparación de tres estrategias de descenso de gradiente}
\end{table}

\begin{knowledgebox}{Beneficios adicionales de los minilotes}
El descenso de gradiente por minilotes no es solo una solución intermedia en términos de eficiencia computacional; el ruido de gradiente que introduce en realidad ayuda al modelo a escapar de mínimos locales pronunciados y tiende a converger hacia regiones de mínimos más planas, las cuales suelen tener mejor capacidad de generalización. Además, el hardware de las GPU está naturalmente diseñado para el cómputo paralelo sobre datos en lote, lo que hace que la velocidad real de entrenamiento del SGD por minilotes supere ampliamente a la actualización muestra por muestra.
\end{knowledgebox}

\subsection{Los retos del descenso de gradiente: mínimos locales y puntos de silla}

En espacios de parámetros de alta dimensión, el landscape de la función de pérdida es sumamente complejo. El descenso de gradiente enfrenta varios retos clave:

\begin{enumerate}
\item \textbf{Mínimos locales} (Local Minima): la función de pérdida puede tener muchos mínimos locales, y el descenso de gradiente puede converger a una posición que no es el óptimo global
\item \textbf{Puntos de silla} (Saddle Points): en espacios de alta dimensión, los puntos de silla son más comunes que los mínimos locales. En un punto de silla, el gradiente es cero en algunas direcciones, lo que estanca la optimización
\item \textbf{Mesetas} (Plateaus): regiones donde el gradiente de la función de pérdida es cercano a cero, lo que hace que las actualizaciones sean extremadamente lentas
\end{enumerate}

\begin{knowledgebox}{Intuición sobre la optimización en espacios de alta dimensión}
En espacios de baja dimensión (como dos dimensiones), los mínimos locales parecen muy comunes. Pero en los espacios de parámetros de millones de dimensiones típicos de las redes neuronales, para que se forme un verdadero mínimo local se necesita que \textbf{todas las dimensiones} se curven hacia arriba al mismo tiempo, lo cual es estadísticamente muy improbable. En realidad, el problema más común en la optimización de alta dimensión son los puntos de silla, en los que algunas dimensiones suben y otras bajan. Esto también explica por qué el aprendizaje profundo funciona tan bien en la práctica: aunque el landscape de la pérdida parezca complejo, la mayoría de los «pozos» tienen una vía de escape.
\end{knowledgebox}

\subsection{Resumen de la sección}

El núcleo del entrenamiento de una red neuronal consiste en minimizar la función de pérdida mediante el descenso de gradiente. La retropropagación utiliza la regla de la cadena para calcular el gradiente de manera eficiente. La tasa de aprendizaje es el hiperparámetro más determinante: tanto un valor demasiado alto como uno demasiado bajo generan problemas, y los algoritmos de tasa de aprendizaje adaptativa (como Adam) pueden ajustarla automáticamente. El SGD por minilotes logra un buen equilibrio entre eficiencia computacional y calidad del gradiente. Aunque la optimización en espacios de alta dimensión enfrenta retos como los mínimos locales y los puntos de silla, los algoritmos de optimización modernos tienen un desempeño sobresaliente en la práctica.
